\section{Ordenação interna e externa}
A ordenação pode ser classificada em dois tipos principais: interna e externa, dependendo da forma como os dados são manipulados durante o processo de ordenação.

Ordenação interna ocorre quando todos os dados a serem ordenados cabem na memória principal do computador (RAM). Esse tipo de ordenação é amplamente utilizado em cenários onde o volume de dados é relativamente pequeno e o acesso rápido à memória permite um desempenho eficiente. Exemplos de algoritmos de ordenação interna incluem Bubble Sort, Merge Sort e Quick Sort.

Ordenação externa, por outro lado, é aplicada quando o volume de dados é tão grande que excede a capacidade da memória principal, necessitando do uso de memória secundária, como discos rígidos ou SSDs. Nesse caso, os algoritmos devem ser projetados para minimizar o número de acessos à memória secundária, que é consideravelmente mais lenta do que a memória principal. Exemplos incluem algoritmos de ordenação por blocos e Merge Sort externo.

Neste trabalho, o foco será exclusivamente na ordenação interna, abordando os principais algoritmos utilizados, suas características e desempenhos em diferentes cenários.